\section{When the agent may choose}\label{sec:selection}

The gate now reads the agent's arguments as \emph{proposals}. A proposed value is bound without the
customer only if the request named it, or an $\Au$ record carries it. So the agent chooses among the
customer's own values, where the gate alone would have required there to be exactly one. This
changes the permission, and with it what memory must keep.

\begin{definition}[Selection permission]\label{def:selection}
$\Perm_{\mathrm{sel}}(h, (a, \theta)) = 1$ if and only if $\Out_a(h) \neq \emptyset$ and:
\begin{enumerate}
\item[(ii')] if $s_a \neq \bot$: $\theta(s_a) \in O$, where $O$ is the set of objects named in
  $\Out_a(h)$ if any is named, and $\Vals_{\kappa_a(s_a)}(h)$ otherwise. Let $\Lambda_\theta$ be the
  requests of $\Out_a(h)$ naming $\theta(s_a)$ or no object; it must be non-empty;
\item[(iii')] for every $p \neq s_a$: $\theta(p)$ is among the values $\Lambda_\theta$ names for $p$ if it
  names any, and among $\Vals_{\kappa_a(p)}(h)$ otherwise.
\end{enumerate}
The definition applies in both regimes.
\end{definition}

$\Perm \le \Perm_{\mathrm{sel}}$: a unique value is in particular a member. What selection gives up
is uniqueness. In the static regime a corrected value stays selectable, so under selection a
correction must be a retraction (Definition~\ref{def:revocable}) to take effect.

\begin{proposition}[The gate implements selection]\label{prop:gate-selection}
If $W$ is faithful and complete (Definition~\ref{def:faithful}; lowered customer messages allowed),
then for every proposal $\pi$, $G(W(\view(h)), \pi)$ executes $o$ only if
$\Perm_{\mathrm{sel}}(h, o) = 1$. Completeness remains necessary: with an unwritten customer
request that names an object, $G$ can execute on an object the request excludes.
\end{proposition}
\begin{proof}
As for Theorem~\ref{thm:sound}, with membership in place of uniqueness.
\begin{itemize}
\item A proposal is bound only when it is named by an $\Au$ request (faithfully an own request) or
  carried by an $\Au$ record (an own value or a books entry).
\item The veto keeps every lowered customer value from being contradicted silently, so a named set
  or object set that $G$ cannot see is never widened.
\end{itemize}
The counterexample to completeness is found by enumeration: the customer asserts object 1, asks for
$a$, then asks for $a$ on object 2. With the last message unwritten, $G$ falls back to object 1,
while $O = \{2\}$.
\end{proof}

\begin{theorem}[What memory must keep when the agent chooses]\label{thm:exact-sel}
Let $\V$ be infinite. In either regime, let the canonical state hold:
\begin{itemize}
\item for each key, the set $\Vals_k$;
\item for each action and each group $g \in \{\bot\} \cup \V$ (the requests naming object $g$, or no
  object), the union of the values its requests name, per parameter.
\end{itemize}
An execution on object $o$ removes the groups $\bot$ and $o$. Reduce by dropping a group $\bot$ that
names nothing when some object group is present. Then $h \equiv h'$ if and only if their reduced
canonical states are equal.
\end{theorem}
\begin{proof}[Proof sketch]
\emph{If.} Permission and updates read only this state, and a bare $\bot$ group beside an object
group adds no name and is spent by every execution.

\emph{Only if.} Membership separates two different sets directly: request a reader naming
everything else, with the missing value proposed. Groups are separated through the object set $O$
and through the named sets, using fresh values as in Theorem~\ref{thm:exact}.
\end{proof}

\begin{corollary}[Choice costs memory]\label{cor:choice}
Letting the agent choose among the customer's values requires remembering every one of them: $N$
bits per key even in the static regime, against $\log_2(N + 2)$ when the gate requires a unique
value (Corollary~\ref{cor:size}). The utility the agent's choice buys is paid for in memory.
\end{corollary}

\begin{theorem}[Safe forgetting when the agent chooses]\label{thm:safe-sel}
In either regime, and up to the equivalence of Theorem~\ref{thm:exact-sel}, a memory state $m$ is
a safe replacement for the truth $t$ if and only if:
\begin{itemize}
\item every key of $m$ holds a subset of $t$'s values;
\item every request group of $m$ is a group of $t$, and names, per parameter, a subset of what $t$'s
  names, non-empty wherever $t$'s is non-empty;
\item every group of $t$ that names an object or a value is a group of $m$. Only a group that names
  neither may be dropped.
\end{itemize}
In words: forget values freely, but never forget what narrows a choice.
\end{theorem}
\begin{proof}[Proof sketch]
\emph{If.} Reads are memberships. A subset permits a subset, and the updates (union, retraction,
expiry, spending, withdrawal) preserve inclusion. A narrowing set kept non-empty never falls back
to the wider key.

\emph{Only if.} Each violation has a continuation.
\begin{itemize}
\item A value $m$ has and $t$ lacks is proposed.
\item A dropped or emptied narrowing set falls back to the key. Assert a fresh value, and $m$
  permits it while $t$ does not.
\item A dropped object group loses its narrowing of the object after the other groups are spent.
  Execute on another object; then a bare request lets $m$ choose an object $t$ excludes.
\end{itemize}
\end{proof}

\paragraph{Two regimes, dual answers.} When the gate needs a unique value
(Theorems~\ref{thm:safe} and~\ref{thm:safe-rev}):
\begin{itemize}
\item forgetting a customer's value is unsafe unless it leaves a mark;
\item a narrowing request may be reduced to a mark.
\end{itemize}
When the agent chooses (Theorem~\ref{thm:safe-sel}):
\begin{itemize}
\item values may be forgotten freely;
\item what narrows a choice may not be forgotten at all;
\item forgetting a retraction (a key larger than the truth) and forgetting a withdrawal (a group
  the truth no longer has) are both unsafe, as before.
\end{itemize}
What a consolidator must preserve therefore depends on how the gate decides, and the gate's design
fixes the consolidation contract (Corollary~\ref{cor:contract}).
